% Generated draft: requires mathematical and source review.
\documentclass[11pt]{article}
\usepackage[margin=25mm]{geometry}
\usepackage{fontspec}
\setmainfont{Latin Modern Roman}
\usepackage{amsmath,amssymb,mathtools,unicode-math}
\setmathfont{Latin Modern Math}
\usepackage{enumitem,needspace,xurl,hyperref}
\setlength{\parindent}{0pt}
\setlength{\parskip}{5pt}
\setlength{\emergencystretch}{3em}
\newcommand{\N}{\mathbb{N}}
\newcommand{\Z}{\mathbb{Z}}
\newcommand{\Q}{\mathbb{Q}}
\newcommand{\R}{\mathbb{R}}
\newcommand{\C}{\mathbb{C}}
\newcommand{\F}{\mathbb{F}}
\newcommand{\A}{\mathbb{A}}
\newcommand{\PP}{\mathbb{P}}
\newcommand{\E}{\mathbb{E}}
\newcommand{\eps}{\varepsilon}
\newcommand{\dd}{\,\mathrm{d}}
\begin{document}
\section*{UnsolvedMath ID: 137}
\textbf{Source problem code:} catalogue entry
\subsection{Definitions and mathematical statement}
Let $n$ be a nonnegative integer and let $A\subset\R^2$ be a finite set of $n$ distinct points. An axis-parallel rectangle is a set $R=[a,b]\times[c,d]$, with real endpoints $a\le b$ and $c\le d$. Call $R$ $1000$-rich for $A$ if $|A\cap R|\ge1000$. A subset $A'\subseteq A$ hits the rich rectangles if
\[
 \forall a\le b\ \forall c\le d,\qquad
 |A\cap([a,b]\times[c,d])|\ge1000
 \ \Longrightarrow\ A'\cap([a,b]\times[c,d])\ne\varnothing.
\]
For every even $n$ and every such $A$, does there exist a hitting subset $A'$ with $|A'|=n/2$? This is the literal half-size formulation in the supplied statement whenever its cardinality is integral. Replacing equality by $|A'|\le n/2$ gives the same question for even $n$, since a smaller hitting subset can be enlarged inside $A$. The record does not specify a rounding convention for odd $n$; an all-$n$ version with $|A'|\le\lfloor n/2\rfloor$ is a possible interpretation, not an established correction of the source wording. Here ``containing $1000$ points'' is understood as containing at least $1000$, in accordance with the supplied background's rich-rectangle formulation. Closed rectangles fix a boundary convention; on a finite point set, the traces of other usual bounded axis-parallel rectangle boundary conventions can be realized by shifting endpoints slightly.

The registry labels the problem open. Its literature triage dated 17 August 2026 reports that no theorem giving a half-size hitting subset for all $1000$-rich rectangles was located in the checked sources, and that no distinct partial result was verified. The requested alternative is a proof of this claim or a counterexample; the dated search assessment is not a proof of present-day openness.
\subsection{Short English statement}
Can half of any finite planar point set be selected so that every axis-parallel rectangle containing at least $1000$ original points contains a selected point? The exact half-size wording is integral for even cardinality; the source leaves odd-cardinality rounding unspecified.
\subsection{Sources}
\begin{itemize}
\item Local UnsolvedMath registry, permanent ID $137$, source code \texttt{catalogue entry}, ``Sparse Hitting Set for Rectangles''; supplied statement and background, attributed to Ben Green's 100 Open Problems.
\item B. Green, \emph{100 Open Problems}, updated January 2026, as identified in the supplied literature triage: \url{https://people.maths.ox.ac.uk/greenbj/papers/open-problems.pdf}.
\item Supplied literature triage checked 17 August 2026: classified open; no half-size theorem or distinct verified partial progress located in its checked source. This draft has not independently consulted the linked source.
\end{itemize}
\end{document}
